% Chapter 9: Agents, Delegation & Parallel Work
\chapter{Agents, Delegation, and Parallel Work}
\label{ch:agents}

\begin{itshape}
A code review needs four perspectives: correctness, security, performance,
and style. Running them sequentially takes 20 minutes.
Running them in parallel takes 5.
But when should you delegate, and when should you just do it yourself?
\end{itshape}

In \cref{ch:extending}, you met subagents briefly as part of the extension
hierarchy. This chapter covers the full delegation system: when to use it,
how to design agent workflows, and how to run work in parallel.

\section{Subagents: Isolated Context Workers}
\label{sec:subagents}

\marginnote[Official]{Define in \code{.claude/agents/}. Each agent runs with its
own context, isolated from the main session.\sourceurl{https://code.claude.com/docs/en/sub-agents}}

\marginnote[Vocab]{In this handbook, \textbf{subagent} means an isolated worker
with its own context. \textbf{Agent team} means a collaborative group
that shares a task list and can message each other.}

\official{Subagents are purpose-built Claude instances that run in isolated context.}
They receive a task, execute it with their own tool set, and return results
to the parent session.

\begin{minted}{markdown}
# .claude/agents/data-quality-reviewer.md
---
name: data-quality-reviewer
description: Review data pipeline for quality issues
tools: [Read, Grep, Glob]
---
You are a data quality specialist.
Review the provided pipeline code for:
1. Schema violations and type mismatches
2. Null handling patterns (silent drops
   vs explicit imputation)
3. Distribution drift between train and serve
4. Data leakage between train and test splits
Report findings with severity ratings
(Critical/High/Medium/Low).
\end{minted}

\margintip{Note the frontmatter key is \code{tools} (not \code{allowed\_tools}).
Skills use \code{allowed-tools} (hyphenated). The inconsistency is intentional---they
follow different specification standards.}

\subsection{Built-In Subagents}

Claude Code ships with several built-in subagent types that can be invoked
via the Task tool without custom definitions:

\begin{description}
  \item[Explore] Fast codebase exploration---find files, search code, answer questions.
  \item[Plan] Read-only analysis for implementation planning.
  \item[Bash] Command execution specialist.
  \item[General-purpose] Multi-tool agent for complex tasks.
\end{description}

% -------------------------------------------------------------------
\section{When to Delegate}
\label{sec:delegation-decision}

\begin{decisionbox}[Delegate or Do Directly?]
\textbf{Delegate when:}
\begin{itemize}
  \item The task is independent of your current work
  \item It requires heavy context that would pollute your main session
  \item Multiple tasks can run in parallel for wall-clock savings
  \item You need an unbiased second opinion (writer/reviewer pattern)
\end{itemize}

\textbf{Do directly when:}
\begin{itemize}
  \item The task takes under 2 minutes
  \item It depends on your current session's context
  \item Sequential reasoning is required (step B needs step A's result)
  \item The task is too simple to justify the overhead
\end{itemize}
\end{decisionbox}

\margincost{Each subagent consumes its own context tokens.
4 parallel agents = 4$\times$ tokens, but only 1$\times$ wall-clock time.
Parallelism trades cost for speed.}

% -------------------------------------------------------------------
\section{Git Worktrees: Isolated File Systems}
\label{sec:worktrees}

\marginnote[Official]{The \code{--worktree} flag creates an isolated copy
of your repository for parallel work.\sourceurl{https://code.claude.com/docs/en/common-workflows}}

Worktrees solve the fundamental problem of parallel file modifications:
two agents editing the same files will conflict. Worktrees give each
agent its own isolated copy of the repository.

\begin{fullwidth}
\begin{beforeafter}
\before
\begin{minted}{bash}
# Manual: create worktree, remember to clean up
git worktree add ../feature-cache -b feature-cache
cd ../feature-cache && claude
# ... work ...
cd - && git worktree remove ../feature-cache
\end{minted}

\after
\begin{minted}{bash}
# Automatic: Claude manages the worktree lifecycle
claude --worktree "Implement caching layer"
claude --worktree "Add rate limiting"
\end{minted}
\end{beforeafter}
\end{fullwidth}

Each worktree gets its own git branch and a complete isolated copy of the
repository. Cleanup is automatic: if no changes are made, the worktree
is removed silently. If changes exist, you are prompted to merge or discard.

\subsection{Subagent Worktrees}

\official{Subagents can also run in isolated worktrees using the
\code{isolation} frontmatter key:}

\begin{minted}{markdown}
# .claude/agents/parallel-reviewer.md
---
name: parallel-reviewer
tools: [Read, Grep, Glob]
isolation: worktree
---
Review the code for correctness and edge cases.
Report findings with file:line references.
\end{minted}

When \code{isolation: worktree} is set, the subagent operates on
its own copy of the repository. This is essential for agents that
modify files---without isolation, two agents editing the same file
will produce conflicts.

\begin{warnbox}[Worktree Housekeeping]
Add \code{.claude/worktrees/} to your \code{.gitignore}.
Claude stores worktree metadata there; it should not be committed.
If orphaned worktrees accumulate (e.g., after a crash), run
\code{git worktree prune} to clean them up.
\end{warnbox}

\margincost{Worktrees trade disk space for isolation.
Each worktree shares git objects but duplicates the working tree.
A 500\,MB repo spawning 4 worktrees may use $\sim$1\,GB additional disk.
Run \code{git worktree prune} to clean up stale entries.}

% -------------------------------------------------------------------
\section{The Writer/Reviewer Pattern}
\label{sec:writer-reviewer}

\official{Use separate sessions to avoid confirmation bias:}

\begin{enumerate}
  \item \textbf{Writer session}: Makes changes based on requirements.
  \item \textbf{Reviewer session}: Reviews the diff without knowledge of intent.
    ``Review this diff for bugs, regressions, and missed edge cases.
    You have no context about why these changes were made.''
\end{enumerate}

This pattern catches errors that the writer is blind to because of
its accumulated context about the ``correct'' approach.

\marginnote[Cross-Ref]{Think together first (\cref{ch:thinking-partner}),
then separate for independent review.}

\practitioner{Extend to multi-perspective review with subagents:}

\begin{enumerate}
  \item \textbf{Correctness agent}: ``Does this logic do what it claims?''
  \item \textbf{Data quality agent}: ``Schema violations? Null patterns? Drift?''
  \item \textbf{Statistical validity agent}: ``Are assumptions met? Distributions stable?''
  \item \textbf{Reproducibility agent}: ``Seeds set? Configs versioned? Results deterministic?''
\end{enumerate}

Running four agents in parallel takes the same wall-clock time as one.

% -------------------------------------------------------------------
\section{Agent Teams (Research Preview)}
\label{sec:agent-teams}

\marginnote[Official]{Research preview. Multiple agents share a task list
and can message each other.\sourceurl{https://code.claude.com/docs/en/agent-teams}}

\official{Agent Teams extend subagents with collaboration:
shared task lists, direct messaging between agents,
and team lead orchestration.}

Best suited for:
\begin{itemize}
  \item Research tasks requiring multiple perspectives
  \item Cross-functional coordination (feature engineering + model training + evaluation + data validation)
  \item Complex multi-module changes that need coordination
\end{itemize}

\practitioner{Prefer parallel subagents for independent tasks that share no state.
Use Agent Teams when agents need to coordinate mid-task---e.g., one agent discovers
a schema change that another agent must respect.}

\begin{warnbox}[Research Preview]
Agent Teams are a research preview feature. The API and behavior
may change. Use for exploration, not production workflows.
\end{warnbox}

% -------------------------------------------------------------------
\section{Quick Reference}

\begin{itemize}
  \item Subagents: isolated context, independent tools, return results to parent
  \item Subagent frontmatter key: \code{tools} (not \code{allowed\_tools})
  \item Worktrees: isolated file systems for parallel file modifications
  \item Writer/Reviewer: separate sessions prevent confirmation bias
  \item Delegate when: independent, context-heavy, parallelizable, needs fresh perspective
  \item Do directly when: quick, context-dependent, sequential, simple
  \item Agent Teams: research preview for multi-agent collaboration
\end{itemize}

% -------------------------------------------------------------------
\begin{keyinsight}
Parallelism in Claude Code is not about speed alone---it is about
\emph{perspective isolation}. A reviewer who cannot see the writer's
reasoning catches different bugs. A security agent that only sees code
(not intent) finds different vulnerabilities. The diversity of perspectives
is the value, not just the time savings.
\end{keyinsight}

% -------------------------------------------------------------------
\begin{trybox}
Identify one review task that currently takes you more than 10 minutes.
Design a 2-agent parallel split:
\begin{enumerate}
  \item Agent 1: data quality review (schema, nulls, distributions).
  \item Agent 2: statistical validity review (leakage, assumptions, reproducibility).
  \item Are the tasks truly independent? (If Agent 2 needs Agent 1's output,
    they cannot run in parallel.)
  \item Write the agent definition file for one of them.
\end{enumerate}
\end{trybox}
